\section{What the Attacker Sees}\label{sec:attacker}

\changed{\Cref{sec:model} built the process. This section classifies what goes in and
out of it, and then hides all but the user processes.}

\subsection{Inputs and outputs}\label{sec:io}

The system receives interrupts. Its input type \code{T_in} is a sum, like the
current pid of \cref{sec:global-state}: \code{inl} is the disk interrupt and
\code{inr} the timer interrupt. \changed{Its output type is the pool's six-tuple
\code{T_out}.}

\begin{definition}[Input and output equivalences]\label{def:io}
The input equivalence $\Ein$ (\m{in_equiv}) and the output equivalence $\Eout$
(\m{out_equiv}) are
\begin{alignat*}{2}
\Ein &\triangleq{} && \priveq + \pubeq, \\
\Eout &\triangleq{} && \hphantom{{}\times{}}(\pubeq + \one^\Lo) \\
&&& {}\times (\priveq + \one^\Lo) \\
&&& {}\times (\pubeq + \one^\Lo) \\
&&& {}\times (\priveq + \one^\Hi) \\
&&& {}\times (\priveq + \one^\Hi) \\
&&& {}\times (\pubeq + \one^\Lo),
\end{alignat*}
one factor per slot, in the order public user, private user, scheduler, NOP,
disk and timer handler, with $\times$ associating to the right, as the output
nests.
\end{definition}

The disk interrupt is thus hidden from the attacker, and the timer interrupt is
public. Of each output the attacker sees whether the slot ran and its value on the
dark slots of \cref{fig:system}, only whether it ran on the light slot, and nothing
on the white ones.
So the attacker sees the U and T steps, and of an S step only that it happened: no
other slot ran. \Cref{sec:gating} shows that the proof forces this shape.

\subsection{The user view}\label{sec:userview}

The \emph{user view} keeps only the two user slots, as one
value of type \code{(Opub + Opriv) + 1}:
\begin{lstlisting}
parse_output o :=
  match o with
  | (Some v, _)         => Some (inl v)
  | (None, (Some v, _)) => Some (inr v)
  | _                   => None
  end
userview m := map id parse_output m
\end{lstlisting}

At $\Lo$ a private-user output and \code{None} are both hidden, so a plain sum
would relate them and hide whether the private user ran. \changed{So the user view
uses $\privvis$ for the private user.}

\begin{definition}[User-view equivalence]\label{def:userview}
The user-view equivalence $\Euser$ (\m{out_equiv_userview}) is
\[
\Euser \triangleq (\pubeq + \privvis) + \one^\emptyset.
\]
\end{definition}

So the attacker still sees when the private user runs. The user view thus erases only public
information: the scheduler and the timer handler, and which of them, or S, ran.

\changed{The private user's input of \cref{sec:loop} pays off here. In this view the
disk interrupt has an effect: on the concrete system of \cref{sec:designs}, the private
user emits \code{Sys} in the run with a disk interrupt where it emits \code{Idle} in
the run without one (\m{trace_with_dI}, \m{trace_no_dI}). The interrupt thus reaches
user space, and only through a value the attacker cannot see.}
